\begin{ex}
Tính: $9 + 7$.
\shortans{16}
\loigiai{
Ta có $9 + 7 = 16$.
}
\end{ex}

\begin{ex}
Tính: $15 - 8$.
\shortans{7}
\loigiai{
Ta có $15 - 8 = 7$.
}
\end{ex}

\begin{ex}
Tính: $8 + 6$.
\shortans{14}
\loigiai{
Ta có $8 + 6 = 14$.
}
\end{ex}

\begin{ex}
Tính: $13 - 9$.
\shortans{4}
\loigiai{
Ta có $13 - 9 = 4$.
}
\end{ex}

\begin{ex}
Tính: $35 + 28$.
\shortans{63}
\loigiai{
Ta có $35 + 28 = 63$.
}
\end{ex}

\begin{ex}
Tính: $72 - 36$.
\shortans{36}
\loigiai{
Ta có $72 - 36 = 36$.
}
\end{ex}

\begin{ex}
Tính: $47 + 35$.
\shortans{82}
\loigiai{
Ta có $47 + 35 = 82$.
}
\end{ex}

\begin{ex}
Tính: $91 - 54$.
\shortans{37}
\loigiai{
Ta có $91 - 54 = 37$.
}
\end{ex}

\begin{ex}
Tính: $60 - 27$.
\shortans{33}
\loigiai{
Ta có $60 - 27 = 33$.
}
\end{ex}

\begin{ex}
Tính: $100 - 45$.
\shortans{55}
\loigiai{
Ta có $100 - 45 = 55$.
}
\end{ex}

\begin{ex}
Tính giá trị biểu thức: $38 + 25 - 19$.
\shortans{44}
\loigiai{
Ta có $38 + 25 = 63$, rồi $63 - 19 = 44$.
}
\end{ex}

\begin{ex}
Tính giá trị biểu thức: $82 - 47 + 29$.
\shortans{64}
\loigiai{
Ta có $82 - 47 = 35$, rồi $35 + 29 = 64$.
}
\end{ex}

\begin{ex}
Điền số thích hợp: $28 + \dots = 61$.
\shortans{33}
\loigiai{
Số cần điền: $61 - 28 = 33$.
}
\end{ex}

\begin{ex}
Điền số thích hợp: $75 - \dots = 38$.
\shortans{37}
\loigiai{
Số cần điền: $75 - 38 = 37$.
}
\end{ex}

\begin{ex}
Điền số thích hợp: $\dots + 36 = 82$.
\shortans{46}
\loigiai{
Số cần điền: $82 - 36 = 46$.
}
\end{ex}

\begin{ex}
Điền số thích hợp: $\dots - 29 = 45$.
\shortans{74}
\loigiai{
Số cần điền: $45 + 29 = 74$.
}
\end{ex}

\begin{ex}
Lớp 2A có 19 bạn nam và 18 bạn nữ. Hỏi lớp 2A có tất cả bao nhiêu bạn học sinh?
\shortans{37}
\loigiai{
Cả lớp có: $19 + 18 = 37$ bạn.
}
\end{ex}

\begin{ex}
Bác nông dân có 64 con gà, bác bán bớt 28 con gà. Bác còn lại bao nhiêu con gà?
\shortans{36}
\loigiai{
Còn lại: $64 - 28 = 36$ con gà.
}
\end{ex}

\begin{ex}
Đoạn dây thứ nhất dài 45 cm, đoạn thứ hai dài hơn đoạn thứ nhất 18 cm. Đoạn dây thứ hai dài bao nhiêu xăng-ti-mét?
\shortans{63}
\loigiai{
Đoạn thứ hai dài: $45 + 18 = 63$ cm.
}
\end{ex}

\begin{ex}
Thùng thứ nhất có 80 lít nước, thùng thứ hai có ít hơn thùng thứ nhất 25 lít nước. Thùng thứ hai có bao nhiêu lít nước?
\shortans{55}
\loigiai{
Thùng thứ hai có: $80 - 25 = 55$ lít.
}
\end{ex}

\begin{ex}
Hai số có tổng bằng 92, số hạng thứ nhất là 47. Số hạng thứ hai là bao nhiêu?
\shortans{45}
\loigiai{
Số hạng thứ hai: $92 - 47 = 45$.
}
\end{ex}

\begin{ex}
Hai số có hiệu bằng 38, số trừ là 25. Số bị trừ là bao nhiêu?
\shortans{63}
\loigiai{
Số bị trừ: $38 + 25 = 63$.
}
\end{ex}

\begin{ex}
Hai số có hiệu bằng 29, số bị trừ là 83. Số trừ là bao nhiêu?
\shortans{54}
\loigiai{
Số trừ: $83 - 29 = 54$.
}
\end{ex}

\begin{ex}
Hiệu của số lớn nhất có hai chữ số và số 35 là bao nhiêu?
\shortans{64}
\loigiai{
Số lớn nhất có 2 chữ số là 99. Hiệu: $99 - 35 = 64$.
}
\end{ex}

\begin{ex}
Tính: $26 + 34 + 15$.
\shortans{75}
\loigiai{
Ta có $26 + 34 = 60$, rồi $60 + 15 = 75$.
}
\end{ex}
